\documentclass[11pt]{article}

\usepackage[margin=1in]{geometry}
\usepackage{amsmath, amssymb}
\usepackage{booktabs}
\usepackage{xcolor}
\usepackage{enumitem}
\usepackage[round]{natbib}
\usepackage{hyperref}

% Color palette (shared with docs/meetings/dfl_explainer.tex)
\definecolor{predBlue}{HTML}{2C6FA8}
\definecolor{optOrange}{HTML}{D97706}
\definecolor{regretRed}{HTML}{B23A48}

\newcommand{\I}{\mathcal{I}}
\newcommand{\B}{\mathcal{B}}
\newcommand{\R}{\mathcal{R}}
\newcommand{\Bi}{\mathcal{B}_i}
\newcommand{\bigM}{M}

\title{A Week-long Discrete Berth Allocation MILP with No-wait Priority
``Services'' and Vessel--Berth Compatibility, for Decision-Focused Learning}
\author{portsDFL --- Port of San Antonio}
\date{\today}

\begin{document}
\maketitle

\begin{abstract}
We formulate the operational problem of planning one week of berth assignments
at the Port of San Antonio as a discrete, dynamic Berth Allocation Problem
(BAP). The model assigns each vessel to a compatible berth and a start time,
minimizing total weighted completion time. A subset of \emph{service} vessels
(the port's term for committed/priority calls) carry a \emph{hard} no-wait
window: they must be berthed essentially on arrival. The service time
$\tau$ that drives the schedule is predicted by an upstream model and the
predictor is trained \emph{decision-focused}, differentiating end-to-end
through the MILP, with a predict-then-optimize baseline as the natural
ablation. This note states the formulation, its hard/soft window modes, the
single-week solve scope, and how it sits in the published BAP literature.
\end{abstract}

\section{Problem setting}
The Port of San Antonio operates several terminals/berths that are
\emph{heterogeneous}: each can physically handle only certain vessel types
(e.g.\ liquid bulk only at the liquid-bulk berth). Over a planning week, vessels
arrive at known times and must each be assigned to a compatible berth and a
start time. Two operational realities drive the formulation:
\begin{itemize}[nosep]
  \item \textbf{Heterogeneous berths.} A vessel may use only the berths
  compatible with its type; the compatibility map is populated data-drivenly
  from the historical $(\text{terminal}\times\text{vessel-type})$ table.
  \item \textbf{Priority ``services''.} A subset of vessels are committed calls
  that must not wait: they are berthed at (or within a small slack of) their
  arrival time. This is a hard reservation, not a soft preference.
\end{itemize}
The service time $\tau_i$ (berth occupation) is not known in advance; it is
predicted from vessel/observed features. The optimizer consumes the prediction
$\hat\tau$, and the predictor is trained so its predictions yield good
\emph{decisions} (\S\ref{sec:dfl}).

\section{Sets, parameters, and variables}
\begin{table}[h]
\centering
\small
\begin{tabular}{@{}ll@{}}
\toprule
Symbol & Meaning \\
\midrule
$\I=\{1,\dots,N\}$ & vessels in the planning week \\
$\B=\{1,\dots,B\}$ & berths \\
$\Bi\subseteq\B$ & berths \emph{compatible} with vessel $i$ \\
$\R\subseteq\I$ & ``service'' (priority, no-wait) vessels \\
$a_i\ge 0$ & arrival time of vessel $i$ (hours from week start) \\
$w_i\ge 0$ & priority weight of vessel $i$ \\
$\tau_i>0$ & service (berth-occupation) time; \emph{predicted} as $\hat\tau_i$ \\
$l_r$ & latest start for service $r\in\R$; $l_r=a_r+\text{slack}$ \\
$\bigM$ & big-$M$ constant ($\ge$ horizon $+$ max berth workload) \\
\midrule
$x_{ib}\in\{0,1\}$ & $1$ if vessel $i$ uses berth $b$ \quad(only for $b\in\Bi$) \\
$s_i\ge a_i$ & start time of vessel $i$ \\
$z_{ijb}\in\{0,1\}$ & $1$ if $i$ precedes $j$ at berth $b$ \\
$\delta_r\ge 0$ & (soft mode) tardiness of service $r$ past its window \\
\bottomrule
\end{tabular}
\caption{Notation.}
\end{table}

\section{Formulation}
The discrete BAP, time-precedence variant
\citep{cordeau2005berth, alzaabi2016berth}:
\begin{align}
\min\quad & \sum_{i\in\I} w_i\,(s_i + \tau_i)
   \;\;\big[\,+\;p\!\!\sum_{r\in\R}\delta_r\ \text{(soft mode)}\,\big]
   \label{eq:obj}\\
\text{s.t.}\quad
& \sum_{b\in\Bi} x_{ib} = 1 && \forall i\in\I \label{eq:assign}\\
& s_i \ge a_i && \forall i\in\I \label{eq:arrival}\\
& s_r \le l_r && \forall r\in\R \ \text{(hard mode)} \label{eq:window}\\
& \delta_r \ge s_r - l_r && \forall r\in\R \ \text{(soft mode)} \label{eq:soft}\\
& z_{ijb} + z_{jib} \le 1 && \forall b,\ i<j:\ i,j\in\Bi\cap\B_j \label{eq:one}\\
& z_{ijb} + z_{jib} \ge x_{ib} + x_{jb} - 1 && \forall b,\ i<j \label{eq:order}\\
& s_j \ge s_i + \tau_i - \bigM\,(1 - z_{ijb}) && \forall (i,j,b),\ i\ne j \label{eq:prec}
\end{align}
\eqref{eq:obj} is total weighted completion time. \eqref{eq:assign} assigns each
vessel to exactly one \emph{compatible} berth---incompatible $x_{ib}$ are never
created (the data-driven analogue of the forbidden-assignment cut
$X_{ij}\le d_{ij}$ of \citealp{alzaabi2016berth}). \eqref{eq:arrival} is the
dynamic arrival release. \eqref{eq:window}/\eqref{eq:soft} encode the no-wait
service window (see \S\ref{sec:windows}). \eqref{eq:one}--\eqref{eq:order} force
an ordering between any two vessels sharing a berth, and \eqref{eq:prec} is the
big-$M$ precedence that prevents overlap. Precedence triples and $z$ variables
are built only for pairs that can share a berth ($b\in\Bi\cap\B_j$), which keeps
the model compact under tight compatibility.

\section{Hard vs.\ soft service windows}\label{sec:windows}
\textbf{Hard mode} (deterministic weekly planner) imposes \eqref{eq:window}:
a service vessel literally cannot start after $l_r$. With $\text{slack}\to 0$
this means ``berth on arrival.'' An over-subscribed week (more concurrent
no-wait services than compatible berths) is then \emph{infeasible}---a genuine
planning signal, surfaced to the planner rather than silently absorbed.

\textbf{Soft mode} (decision-focused training) drops \eqref{eq:window} and adds
the tardiness variable and penalty $\eqref{eq:soft}$, $+\,p\sum_r\delta_r$ in
\eqref{eq:obj}. This keeps every solve \emph{feasible} under any predicted
$\hat\tau$---essential because the DFL trainer re-solves the MILP for arbitrary
(early-training, possibly poor) predictions, and an infeasible solve would abort
training and undefine the regret.

\section{Single-week solve scope}
The MILP solves exactly one 7-day instance ($N$ fixed, horizon $168$\,h).
Selecting the week's vessels is a \emph{pre-solve slicing step}; there is no
multi-week, rolling-horizon, or cross-week carry-over inside the optimizer.
Vessels arriving outside the chosen window are filtered out before the model is
built.

\section{Decision-focused learning integration}\label{sec:dfl}
$\hat\tau$ enters both the objective \eqref{eq:obj} and the precedence
constraints \eqref{eq:prec}. Because the predicted quantity appears in the
constraints (not only as a linear objective cost vector), SPO+
\citep{elmachtoub2022spo} does not apply directly; we differentiate through the
solver with PyEPO blackbox differentiation \citep{pogancic2020blackbox}, which
trains the predictor to minimize downstream regret. A predict-then-optimize
(MSE) predictor is the controlled baseline. Regret is measured against the
full-information schedule (the MILP solved under true $\tau$), and is
$\ge 0$ by construction.

\section{Related formulations}

\paragraph{Related formulations.}
Our model is a discrete, dynamic Berth Allocation Problem (BAP) over a single
planning week, solving an assignment, start-time and per-berth precedence MILP
that minimizes total completion time $\sum_i (s_i + \tau_i)$ (equivalently,
total waiting; vessel priority weights $w_i$ are a planned extension --- the
implemented objective is unweighted). We
situate it against the closest published BAP formulations along five axes: time
representation, objective, time-window/priority semantics, vessel--berth
compatibility, and whether the service (handling) time is \emph{learned}.

The discrete--dynamic backbone descends directly from \citet{imai2001dynamic}
and \citet{cordeau2005berth}. \citet{imai2001dynamic} introduced the discrete
dynamic BAP as a three-index assignment of each vessel to a (berth,
service-order) slot, minimizing the total \emph{unweighted} service time
(waiting plus handling) with berth-dependent handling times $C_{ij}$ and hard
arrival-release constraints, but no per-vessel weights, no due dates, and no
eligibility filter. \citet{cordeau2005berth} recast the same discrete dynamic
problem as a multi-depot vehicle-routing problem with time windows, replacing
service-order indices with arc/precedence variables and, crucially for us,
introducing a \emph{weighted} service-time objective with hard arrival and
due-date windows and berth-dependent handling. Our objective
$\sum_i w_i (s_i + \tau_i)$ is precisely the weighted-completion-time family of
\citet{cordeau2005berth}, and our per-berth precedence variables $z_{ijb}$ with
big-$M$ linking are the time-precedence analogue of their routing arcs; we
differ in using an explicit calendar-time formulation over $s_i$ and $\tau_i$
rather than sequencing arcs, in adding a hard no-wait priority class, and in
treating $\tau$ as a learned quantity.

On the priority axis, \citet{alzaabi2016berth} is the closest single antecedent
to our deterministic core: they minimize $\sum_j w_j (S_j + p_j - a_j)$ -- the
same weighted-completion-time objective net of arrival -- over a discrete
assignment-plus-sequencing MILP with hard arrival release $S_j \ge a_j$. They
also provide the cleanest precedent for our compatibility treatment: an explicit
per-pair feasibility matrix $d_{ij}$ enforced by the forbidden-assignment cut
$X_{ij} \le d_{ij}$ (plus a length/draft capacity variant), which is
structurally identical to our data-driven berth-type eligibility restriction of
$x_{ib}$ to compatible $(i,b)$ pairs. Their priority, however, is purely a soft
cost weight $w_j$; there is no constraint forcing a privileged vessel to start
without waiting. \citet{Ursavas2022PriorityControlBAP} likewise treats priority
softly, but via a dynamic, time-aging priority index and user preference weights
tuned by a scatter-search metaheuristic over a stochastic discrete-event
simulation of heterogeneous berths; preferential vessels are favored in
expectation but never guaranteed. In contrast, our ``service'' vessels (the
port's term for committed/priority calls) receive a \emph{hard} no-wait window
$s_r \le a_r + \mathrm{slack}$ with $\mathrm{slack}\!\to\!0$, an absolute
reservation rather than an objective weight or a dispatch heuristic.

On the time-window axis, \citet{golias2007berth} is the canonical soft-deadline
berth model: each vessel has a requested departure window $[t_{j1}, t_{j2}]$,
with early, timely, and late departures rewarded or penalized economically
through big-$M$ segment-classification constraints, so deadlines are penalized
but never enforced. This is the opposite design choice to ours:
\citet{golias2007berth} deliberately makes windows soft and avoids the hard
early-completion constraint (which it notes can be infeasible), whereas our
priority calls impose exactly such a hard no-wait constraint, and our objective
is weighted completion time rather than an earliness--tardiness net benefit.

The closest work on the learning axis is \citet{chu2025integrating} (with
companion \citealp{chu2026coordinated}), which embeds a vessel-arrival-time
forecast into a continuous BAP using an XGBoost predictor and a real Hong Kong
case study. Critically, theirs is a two-stage \emph{predict-then-optimize}
pipeline: the predictor is trained on a forecasting loss (MAE/RMSE) and its
point forecasts are passed as fixed inputs to the optimizer, with no gradient
flowing back from decisions to predictions. The uncertain quantity is the
\emph{arrival} time, the model is a continuous rectangle-packing BAP, and there
is no priority class or learned berth-type compatibility. Our work instead
trains the service-time predictor $\tau$ \emph{decision-focused}, differentiating
end-to-end through the MILP with PyEPO blackbox differentiation
\citep{pogancic2020blackbox} so the predictor minimizes downstream decision
regret rather than prediction error; we report a predict-then-optimize MSE
baseline (in the spirit of \citealp{elmachtoub2022spo}) as the natural ablation.
The only neighboring port study that is genuinely decision-focused,
\citet{pu2025predictthenoptimize}, applies DFL to seaport
power-logistics/shore-power scheduling via a differentiable convex surrogate with
continual learning, not to the BAP, and models neither priority vessels,
berth-type compatibility, nor weighted completion time. To our knowledge no
prior BAP formulation combines all four of our distinguishing ingredients --
discrete weekly horizon, hard no-wait priority calls, data-driven berth-type
compatibility, and a decision-focused (blackbox-differentiated) service-time
predictor.

\begin{table}[t]
\centering
\small
\setlength{\tabcolsep}{4pt}
\renewcommand{\arraystretch}{1.2}
\begin{tabular}{@{}p{2.7cm} p{2.3cm} p{2.4cm} p{2.2cm} p{2.0cm} p{2.3cm} p{1.6cm}@{}}
\toprule
\textbf{Paper} & \textbf{Model type} & \textbf{Objective} & \textbf{Time windows} & \textbf{Priority} & \textbf{Berth compat.} & \textbf{Pred./DFL} \\
\midrule
\citet{imai2001dynamic}            & Discrete, service-order MIP        & Total service time (unweighted)        & Arrival hard; no due date        & None                         & Berth-dep.\ handling $C_{ij}$         & No / No \\
\citet{cordeau2005berth}           & Discrete, MDVRPTW (arc/prec.)      & Weighted service time                  & Arrival + due date, both hard    & Soft weight $v_i$            & Berth-dep.\ handling $t_{ik}$         & No / No \\
\citet{golias2007berth}            & Discrete, position-indexed MILP    & Earliness--tardiness net benefit       & Departure window, \emph{soft}    & Soft monetary premia         & Berth-dep.\ handling $C_{ij}$         & No / No \\
\citet{alzaabi2016berth}           & Discrete, big-$M$ sequencing MILP  & Weighted completion $\sum w_j(S_j{+}p_j{-}a_j)$ & Arrival hard; no hard due date & Soft weight $w_j$         & Explicit matrix $X_{ij}\le d_{ij}$    & No / No \\
\citet{Ursavas2022PriorityControlBAP} & Sim.--opt.\ (scatter search)    & Expected total service time            & Stochastic ETA, soft             & Soft aging index + weights   & Per-class rates + thresholds          & No / No \\
\citet{chu2025integrating}         & Continuous BAP (rectangle pack)    & Delay-induced cost / turnaround        & Pred.\ arrival release           & None                         & Length/tide/position (physical)       & Yes (arrival) / No (PtO) \\
\citet{pu2025predictthenoptimize}  & Convex surrogate (not BAP)         & Power-logistics scheduling cost        & ---                              & None                         & --- (not berths)                      & Yes / Yes (convex DFL) \\
\midrule
\textbf{Ours}                      & \textbf{Discrete, weekly, prec.\ MILP} & \textbf{Weighted completion $\sum w_i(s_i{+}\tau_i)$} & \textbf{Arrival hard; no-wait $s_r\le a_r$ hard} & \textbf{Hard no-wait class + weights} & \textbf{Data-driven type matrix} & \textbf{Yes ($\tau$) / Yes (blackbox DFL)} \\
\bottomrule
\end{tabular}
\caption{Closest published BAP formulations versus ours. ``Berth-dep.\
handling'' = heterogeneity captured only through the berth-dependent
handling-time matrix (no explicit eligibility filter); ``explicit matrix'' /
``data-driven type matrix'' = a hard vessel--berth feasibility restriction.
``PtO'' = two-stage predict-then-optimize (predictor trained on a forecasting
loss); ``DFL'' = decision-focused, predictor differentiated through the
optimizer.}
\label{tab:related_bap}
\end{table}

\paragraph{Novelty.}
What is genuinely novel in ours is the \emph{combination}, within a single
discrete weekly BAP, of (i) a \emph{hard} no-wait reservation for a committed
``service''-vessel class ($s_r \le a_r + \mathrm{slack}$,
$\mathrm{slack}\!\to\!0$) and (ii) a \emph{decision-focused} (PyEPO
blackbox-differentiated) predictor of the \emph{service/handling} time $\tau$
that enters both the objective and the precedence constraints. No prior BAP
formulation we found pairs a hard no-wait priority constraint with an
end-to-end-differentiated service-time predictor: the priority-focused works
\citep{Ursavas2022PriorityControlBAP, alzaabi2016berth, cordeau2005berth,
golias2007berth} all treat priority \emph{softly} (objective weights, monetary
premia, or an aging dispatch index) and treat handling time as a known input;
the only closely related learning-based port work \citep{chu2025integrating} is
two-stage predict-then-optimize on the \emph{arrival} time with no
decision-aware training, no priority class, and a continuous (not discrete
weekly) quay; and the only genuinely decision-focused port study
\citep{pu2025predictthenoptimize} targets power-logistics, not berth allocation.
A secondary, more incremental contribution is that the predicted quantity is the
service time entering the precedence constraints (so SPO+ is inapplicable and a
blackbox/perturbation-based gradient is required), a different and somewhat
harder DFL setting than the objective-only cost-vector prediction used in most
predict-then-optimize benchmarks.

What is \emph{not} novel and should be stated plainly: the underlying
optimization model is a standard discrete time-precedence DBAP -- the
assignment, start-time, and per-berth big-$M$ precedence structure and the
weighted-completion-time objective $\sum_i w_i(s_i+\tau_i)$ are the classical
\citet{cordeau2005berth} / \citet{alzaabi2016berth} family, not a new model. The
explicit vessel--berth compatibility encoding ($x_{ib}$ restricted to feasible
pairs) is structurally the forbidden-assignment cut $X_{ij}\le d_{ij}$ of
\citet{alzaabi2016berth}; calling it ``data-driven'' refers only to how the
compatibility matrix is populated (from a vessel-type $\times$ terminal table),
not to a new constraint family. Per-vessel priority weights $w_i$ are likewise
standard. Blackbox differentiation of combinatorial solvers
\citep{pogancic2020blackbox} and the predict-then-optimize / SPO+ framing
\citep{elmachtoub2022spo} are existing methods we apply, not invent. The novelty
is therefore in the problem framing and the integration -- hard committed-call
reservations plus decision-focused service-time learning in a real (San Antonio)
weekly berth-planning setting -- rather than in any single new constraint or new
learning algorithm.


\section{Implementation pointers}
The MILP lives in \texttt{optimizers/src/bap\_optim/discrete\_bap.py}
(\texttt{DiscreteBAP}, with the \texttt{hard\_windows} flag). The instance
descriptor \texttt{BAPInstance} (with optional \texttt{latest\_start},
\texttt{berth\_compat}, \texttt{service}) is in \texttt{bap\_optim/instance.py}; the
berth catalog and compatibility in \texttt{bap\_optim/berths.py}; the pre-solve
weekly slicing in \texttt{bap\_optim/weekly\_instance.py}; and the deterministic
planner in \texttt{prediction\_models/scripts/plan\_week.py}. Open modeling decisions
and assumptions are tracked in \texttt{prediction\_models/docs/formulation/decisions\_and\_questions.md}.

\bibliographystyle{plainnat}
\bibliography{references}

\end{document}
